\section{Reproducibility}
\label{app:reproducibility}

\subsection{Report inputs}

The main training and performance tables and figures are generated from
checked-in data snapshots. The publication data directory contains the
training-loss series, same-accelerator performance aggregates, checkpoint
inventories, evaluation results, plotting programs, and a SHA-256 ledger.
Figure generators verify their declared inputs before writing an output. The
machine-readable recipe ledger records the route, precision policy, rounding
and preconditioning choices, batch geometry, seed when recovered, checkpoint
horizon, and evidence status for every experiment discussed in the paper.
Shared model, data, and optimizer fields are stored separately so that an
incomplete historical row cannot silently inherit a value from a newer run.

For display consistency, the historical BF16 and TE-native training-loss
curves in Figure~\ref{fig:llama8b-loss} use fixed offsets inferred from 41
matched steps in seed-42 controls. The raw series and adjustment remain in the
publication data.

The long-run curves combine only explicitly identified execution segments.
The \localcta{}+fixed-H16 and operand-hybrid trajectories each cross a
documented full-state resume boundary. Their manifests identify the source and
continuation ranges; no overlapping step is counted twice. Training-loss
observations for custom global NVFP4 are regular through step 19,000. The
configured Weights \& Biases run is no longer returned by the service, and the
dense local CSV named by the historical manifest was not present in any
recovered or released tree. Seventy-four exact observations survive after
step 19,000 in fragmented training logs. The largest missing sequence of
expected ten-step records runs from step 24,710 through step 33,000. The
dashed late line is a visual LOWESS fit through the surviving observations,
evaluated at the BF16 plotting steps. It is never used to compute a reported numeric
endpoint, table entry, or statistic.

\subsection{Checkpoint and evaluator contract}

Training comparisons use exact shared steps. Custom-route values are raw
observations. Downstream comparisons use literal step-38,000
checkpoints; a terminal, nearby, or checkpoint from a different training run
is never substituted. Distributed checkpoints must contain all expected
shards and pass conversion checks before evaluation.

The evaluator must reproduce the training model's RoPE and RMSNorm semantics.
The Llama-3.1 8B training configuration uses base $500{,}000$ and the
long-context RoPE scaling parameters $(8,1,4,8192)$. An earlier evaluator
disabled this scaling in both the converted-checkpoint evaluator and native
TorchTitan reference. Agreement between those two implementations therefore
did not establish agreement with the training model. Results from that
evaluator are excluded.

An earlier validation stream used the correct model semantics but failed a
data-source robustness check: apparent route ordering changed when examples
came from different contiguous regions of the source datasets. Those values
are also excluded. The replacement stream was constructed afresh from
deterministic windows in the 90--95\% physical sample range of the DCLM and
OLMo-without-DCLM components. Documents are deduplicated by raw-text SHA-256,
combined at an 82/18 token ratio, globally hash-shuffled, and only then packed
and sharded. Thus no GPU worker is restricted to one source or contiguous
region of source data.

The replacement fixed validation stream contains 768 sequences of 8,193 tokens, yielding
6,291,456 scored next-token targets. It is tokenized once with the pinned
Llama-3.1 tokenizer and reused byte-for-byte for every checkpoint. We call it
a \emph{fixed external validation stream}, not a held-out split, because zero
document overlap across every historical training trajectory has not been
proven. Its source ledger contains 4,597 documents with
4,597 unique raw-text hashes. Before reporting route comparisons, the
corrected checkpoint evaluator must agree exactly with the native TorchTitan
model at the full 8,192-token context, and a late BF16 checkpoint must improve
over an early BF16 checkpoint by more than five sequence-level standard
errors. Both checks pass. The converted checkpoint and native TorchTitan model
agree exactly at full context for the step-2,000 and step-29,000 BF16
checkpoints. Their validation NLL values are 3.283272 and 2.460520,
respectively; the paired late-minus-early difference is $-0.822753$ with
sequence-level standard error 0.003708.

Downstream evaluation uses BF16 inference and one evaluator configuration for
all routes. MMLU is 5-shot, HellaSwag 10-shot, WinoGrande 5-shot, and
ARC-Challenge 25-shot; HellaSwag and ARC report normalized accuracy. A route
enters the common panel only after the exact checkpoint and corrected model
semantics are verified.

\subsection{Historical execution identity}

Historical experiments are bound to the code, submodule, kernel asset,
checkpoint, and output identities recovered for that execution. The current
source and runtime tips are not assigned retroactively to an older run. Where
a field was not recovered, the ledger records it as unknown rather than
guessing the nearest commit. Complete identifiers and content digests are in
the publication data ledger rather than the narrative.

The custom global NVFP4 path uses one orientation-consistent 2D16 weight
encoding. The code used by the completed \mxfp{}+row-SR+fixed-H32 run is
content-verified in its run manifest. The \localcta{}+fixed-H16 and
operand-wise plain-H16 trajectories each join verified execution segments at
an exact full-state checkpoint. Their manifests record implementation changes,
segment boundaries, and checkpoint state.

\subsection{Controlled Dgrad diagnostic}
\label{app:dgrad-diagnostic}

The operand-format diagnostic uses tensors from the step-28,000 checkpoint of
the completed \mxfp{}+row-SR+fixed-H32 trajectory. It computes Dgrad with
\mxfp{} and \localcta{} against the corresponding BF16 result. The diagnostic
configures a paired plain-H16 transform for Wgrad only; that transform does
not enter the scored Dgrad contraction and is not part of the comparison.

For each captured tensor, four stochastic candidate tensors are averaged
before computing
\begin{equation}
  E=\sum_i\left(\widehat{dX}_i-dX_i^{\mathrm{BF16}}\right)^2.
\end{equation}
The capture covers three
sequence-2,048 batches, QKV, attention-output, and FFN projections in all 32
blocks, and an independently seeded stochastic-rounding generator rather than
the training run's saved generator state. This gives 96 layer--component
comparisons per batch and 288 in total.
\localcta{} has lower error energy in all 288 comparisons. Giving each tensor
equal weight, the cross-capture geometric mean of
$E_{\mathrm{MX}}/E_{\mathrm{localCTA}}$ is 1.292.

\subsection{Excluded experiments}

The custom global \nvfp{}/\mxfp{} depth hybrid remains outside the common endpoint panels because
it diverged and retained only 16 of 32 terminal shards. The TE route with four
final BF16 blocks and the fixed-H32 operand hybrid both completed 160B-token
trajectories with resumable final checkpoints and are included in the endpoint
panels.

Diagnostic screens are also excluded from the common panels. The
low-precision output-head configurations produced no stable model-level
speedup. A \localcta{} transform screen resumed near the end of training under
a changed runtime, so it is not a causal comparison. The superseded downstream
evaluator omitted training-time RoPE scaling, while the superseded validation
stream was sensitive to the contiguous source-data region selected. Their
records are retained for audit, not used as scientific endpoints.

Run \texttt{make} in the report directory to regenerate the figures and PDF.
The checked-in data README lists each source file, generator, and verification
command; \texttt{data/SHA256SUMS} binds the publication inputs.

\subsection{Code and data availability}
\label{app:code-release}

Code, pinned kernel sources, experiment recipes, evaluation tools, and figure
builders are available at
\url{https://github.com/MrHuff/mfu-fp4-training}. Project-authored software is
released under the MIT License. The original manuscript and figures are
released under CC BY 4.0; included third-party components retain their own
terms. Training corpora, tokenizer assets, model checkpoints, and checkpoint
metadata are not redistributed. The repository instead provides hash-bound
input schemas and public recipes for binding authorized local copies without
recording credentials or private storage locations.
